\section{问题二：风险修正的日前购电策略}

\subsection{为什么不能直接用``平均预测''}

问题二的短缺电量需按交易时刻电价的 5 倍紧急购电，而计划多买的边际损失仅约为正常电价。因此预测误差的上、下侧成本明显不对称：仅最小化均方误差的点预测并不能直接最小化购电总成本。我们采用``因果点预测 + 残差风险分位''的两层结构。

\subsection{负载与光伏的逐日递推预测}

对第 $d$ 天，在 0:00 时只能使用此前数据。负载的候选模型包括 D-1、D-7、移动均值/中位数等；确认实验中 7 d 季节朴素的走步 MAE 最低，故取
\[
\hat L_{d,t}=L_{d-7,t}.
\tag{5}
\]
光伏使用最近 7 d 指数加权均值
\[
\hat P_{d,t}=\sum_{k=1}^{7}w_kP_{d-k,t},\quad \sum w_k=1,
\tag{6}
\]
最近日期权重更大。2--12 月负载、光伏 MAE 分别为 176.796 kW 和 149.774 kW，组合净负荷 MAE 为 271.185 kW。

\subsection{0.8 分位风险储备}

定义净负荷预测残差
\[
r_{d,t}=(L_{d,t}-P_{d,t})-(\hat L_{d,t}-\hat P_{d,t}).
\tag{7}
\]
制定第 $d$ 天计划时，只使用过去至多 28 d 的同时间隔残差，取 0.8 分位数：
\[
R_{d,t}=Q_{0.8}\{r_{j,t}:d-28\le j<d\}.
\tag{8}
\]
再令 $\tilde N_{d,t}=\hat L_{d,t}-\hat P_{d,t}+R_{d,t}$ 进入式（1）。0.8 的经济解释来自 newsvendor 临界分位：计划多买 1 kWh 的过量成本约为 $p$；计划少买 1 kWh 后通过紧急购电补足，其相对常规购电的增量损失约为 $4p$，因此临界分位为 $4/(4+1)=0.8$。

实际净负荷实现后，若计划购电与储能供电仍不足，则
\[
e_{d,t}=\max\{(L_{d,t}-P_{d,t})\Delta t-(g_{d,t}+d_{d,t}-c_{d,t}),0\},
\tag{9}
\]
全天费用为
\[
C_{2,d}=\sum_t p_tg_{d,t}+\sum_t5p_te_{d,t}.
\tag{10}
\]

\subsection{全年结果与对照方案}

2025-02-01 至 2025-12-31 的计划购电量为 22,068,285.289 kWh，紧急购电量为 367,486.925 kWh；计划购电费 13,496,102.516 元，紧急购电费 1,419,991.140 元，总费用 \textbf{14,916,093.656 元}。与``无储能但采用相同 0.8 分位风险计划''的约 18,609,128.71 元相比，下降约 \textbf{19.85\%}，说明储能带来的跨时段配置收益并非由复杂预测模型伪造。
